%Chapter 4
\section{Methodology}
\label{sec:methodology}
Chapter~\ref{sec:design} established the design of the two source-code versions and documented the requirements constraining the claims made. This chapter specifies the procedures used to generate evidence supporting those claims. Three protocol artefacts are evaluated. The first is the PQXDH baseline, defined as the unmodified upstream libsignal implementation from which the project fork originates. It is pinned at release v0.73.3 (commit \texttt{7cce36e9}), which was the final release before SPQR was integrated into libsignal \cite{noauthor_release_nodate,noauthor_release_nodate-2}. Version 1 pairs ML-KEM-1024 with ML-DSA-87 within the fully post-quantum handshake scope defined in Section~\ref{subsubsec:req-functional}. Version 2 substitutes HQC-256 for ML-KEM-1024 within the same handshake structure described in Section~\ref{subsubsec:variant-spec}. Each value reported in Chapter~\ref{sec:results} is associated with a specific, pinned artefact measured under a documented procedure rather than interpreted as a general claim about an algorithm.

This chapter follows the reporting rule established at the end of Chapter~\ref{sec:design}. Under this rule, Chapter~\ref{sec:results} distinguishes demonstrated findings, conditionally supported claims, unresolved issues, and unmet requirements. A conditional claim identifies the condition on which it depends. An unresolved issue lacks sufficient evidence. An unmet requirement is one that inspection shows the artefact does not satisfy.

%Chapter 4.1
\subsection{Artefacts, Benchmark Environments, and Reproducibility}
\label{subsec:environment}

All reported measurements in this project pertain to specific, pinned builds of the code rather than to algorithms in the abstract. Protocol integration and benchmark entry points are restricted to the Rust \texttt{libsignal-protocol} crate within a fork of libsignal~\cite{farrugia_marksaviourpost-quantum-signal-v7_2026}. Version 2 also incorporates the patched \texttt{hqc-kem} dependency described below, which sits outside the Cargo workspace. The Java, Swift, and Node.js bindings are out of scope, and all builds target the Rust crate directly. The PQXDH baseline corresponds to project commit \texttt{26ff061} (tag \texttt{v0.0.0}), which imports upstream release v0.73.3. This release was selected because it is the final version before SPQR integration. It therefore enables measurement of the PQXDH-to-Double-Ratchet path without the subsequent post-quantum ratchet. At the protocol level, it retains the published hybrid PQXDH construction used by the deployed comparator. Its pinned handshake path uses Kyber-1024. It is a reproducible pre-SPQR source configuration, not a byte-for-byte reproduction of current production clients (Section~\ref{subsubsec:req-nonfunctional}); the import was verified tree-identical to upstream outside the measured crate, the only differences being continuous-integration and tooling files. Commit \texttt{f55ae91} (tag \texttt{v1.0.0}) converts that baseline to the fully post-quantum form described in Chapter~\ref{sec:design}. Version 2, tagged \texttt{v2.0.0} at commit \texttt{95f53fac}, is derived from Version 1 in a single commit and includes the vendored dependency described below. The baseline-to-Version-1 comparison therefore measures the aggregate cost of the fully post-quantum transition. The Version-1-to-Version-2 comparison maintains the handshake structure, ML-DSA-87 authentication, and measurement procedure while replacing the KEM suite. It measures the complete substitution within the pinned implementations, but it does not isolate algorithmic cost from differences in library maturity or optimisation.

Measurement harnesses are maintained on separate immutable checkouts so the protocol tags remain unchanged: a baseline checkout from \texttt{26ff061}, session-harness checkouts from \texttt{v1.0.0} and \texttt{v2.0.0}, and a common primitive-harness checkout from \texttt{v2.0.0} with both KEM features enabled. Each diff is limited to benchmark sources, and these checkouts are not additional protocol variants. Section~\ref{subsec:perf-measurement} specifies their targets and equivalence checks.

Dependency versions are fixed across all measurements: ML-DSA-87 is \texttt{libcrux-ml-dsa} 0.0.8 \cite{noauthor_cratesio_2026}, and HQC-256 is \texttt{hqc-kem} 0.1.0, vendored at \texttt{third-party/hqc-kem} with a correction to an inverted length check in its byte-deserialisation macro. The committed \texttt{Cargo.lock} records the complete transitive set, and the \texttt{mlkem1024} and \texttt{hqc256} features select the KEMs at build time. Appendix~\ref{app:dependency-detail} records the full vendoring rationale and feature configuration; the resulting threat to validity is addressed in Section~\ref{subsec:threats-validity}.

All artefacts are built and measured on two machines. The primary machine, which produces every headline result in Chapter~\ref{sec:results}, is a Ryzen 9 9950X3D with 48 GB of DDR5-8400 memory running Windows 11 and targeting \texttt{x86\_64-pc-windows-msvc}. The build also requires the Microsoft Visual C++ (MSVC) toolchain and the Protocol Buffers compiler \texttt{protoc} 35.1, invoked by the crate's build script via \texttt{prost-build}. The secondary machine is an Apple MacBook Pro with an M1 Pro chip and 16 GB of LPDDR5-6400 unified memory, running macOS and targeting \texttt{aarch64-apple-darwin}.

Together the pinned commits, lockfile, vendored dependencies, feature flags, and specified commands support replication on compatible systems. Serialised lengths are exactly reproducible for the fixed inputs of Section~\ref{subsec:size-measurement}; timings are reproducible across comparable hardware, with claims resting on relative orderings. Randomness is drawn from each operating system's CSPRNG, and all raw outputs, commands, harness diffs, and manifests are retained in the project evidence archive.

%Chapter 4.2
\subsection{Correctness Verification}
\label{subsec:correctness}

Before any timing measurements are accepted, each pinned artefact undergoes executable functional checks. These checks demonstrate only the exercised behaviours described below. They do not establish protocol correctness or replace the requirements, scope qualifications, and recorded implementation limitations of Chapter~\ref{sec:design}. Version 1 is tested with \texttt{cargo test -p libsignal-protocol --no-default-features --features mlkem1024}, and Version 2 with \texttt{cargo test -p libsignal-protocol --no-default-features --features hqc256}. \texttt{test\_full\_pq\_prekey\_handshake} and \texttt{test\_alice\_and\_bob\_agree\_with\_one\_time\_kem\_prekey} establish agreement in last-resort-plus-one-time mode. \texttt{test\_prekey\_handshake\_without\_one\_time\_kem} and \texttt{test\_alice\_and\_bob\_agree\_without\_one\_time\_kem\_prekey} establish last-resort-only mode, thereby demonstrating both implemented modes of Section~\ref{subsubsec:req-functional}.

\texttt{test\_bad\_signed\_pre\_key\_signature} and \texttt{test\_bad\_kyber\_pre\_key\_signature} alter the signed-ratchet-key and last-resort-KEM signatures respectively and require rejection, while \texttt{test\_bob\_rejects\_bad\_transcript\_signature} does the same for the initiator's transcript signature. \texttt{prekey\_message\_failed\_decryption\_does\_not\_update\_stores} demonstrates rejection of a tampered authenticated inner message without persisting a session. \texttt{test\_repeat\_bundle\_message} verifies that two separately encrypted \texttt{PreKeySignalMessage}s from one unacknowledged outbound session can both be processed against the responder's established session. It does not replay identical message bytes or test duplicate-delivery suppression. Adapter tests exercise key generation, encoded-length checks, encapsulation, and decapsulation for each compiled KEM, and reject a wrong-type ciphertext offered to an HQC-256 key. Version 2 exercises the corrected HQC byte constructors on valid-length encapsulation keys, decapsulation keys, and ciphertexts through the adapter's encapsulation--decapsulation round trip. Neither the adapter nor the vendored tests probes adjacent invalid lengths. The vendored crate's known-answer and round-trip tests are run separately with \texttt{cargo test --manifest-path third-party/hqc-kem/Cargo.toml}. The example \texttt{examples/full\_session.rs} provides a runnable public-API demonstration but does not substitute for these assertions. The baseline is tested independently at project commit \texttt{26ff061} with \texttt{cargo test -p libsignal-protocol}, and its pass/fail counts are recorded before any baseline measurements are accepted.

Coverage remains narrower than the negative-test requirement of Section~\ref{subsubsec:req-functional}. No protocol-level test separately alters the optional one-time-KEM signature, identity or suite identifiers, individual authenticated transcript fields, signed public keys as distinct from their signatures, KEM ciphertext fields, or malformed combinations of the added fields. Adjacent invalid HQC key and ciphertext lengths are likewise not tested directly against the corrected vendored constructors. No test passes one version's encodings to the other, and no end-to-end test replays an identical initial-message byte sequence. The absent test requirements are classified as unmet. The corresponding protocol properties remain unresolved unless a separately stated condition supports a conditional result.

Chapter~\ref{sec:results} records a requirement-to-evidence matrix for the principal functional, security, and non-functional claims evaluated directly. Detailed criteria remain in Appendix~\ref{app:extended-requirements}; the matrix is a compact claims summary rather than an exhaustive checklist. This prevents a passing functional test from being treated as evidence for an untested security property.

%Chapter 4.3
\subsection{Performance Measurement}
\label{subsec:perf-measurement}

Timing data are collected using Criterion 0.5.1, as specified in the committed lockfile. Unless an invocation explicitly records an override, each case uses Criterion's default settings. These consist of 100 samples, a three-second warm-up, and a five-second target measurement time \cite{heisler_bheislercriterionrs_2026}. The median point estimate and its 95\% bootstrap confidence interval are reported. Criterion's Tukey outlier classifications and the complete sampled distribution are retained. Classified outliers are not excluded from the estimates.

The corrected KEM target (\texttt{benches/kem.rs}) measures \texttt{generate}, \texttt{encapsulate}, and \texttt{decapsulate} separately for HQC-256, ML-KEM-1024, and Kyber-1024, all at NIST category~5. Ciphertext--secret-key fixtures for decapsulation are fully materialised before the timed closure so that encapsulation is not included in the decapsulation result. The target is adapted from the upstream benchmark by replacing the Kyber-768 comparison with HQC-256 and ML-KEM-1024 and by adding raw key, ciphertext, and shared-secret length reporting. Type-prefixed protocol encodings are reported separately in Section~\ref{subsec:size-measurement}. The target runs from the common primitive-harness checkout derived from \texttt{v2.0.0} with \texttt{hqc256} and \texttt{mlkem1024} enabled. This is justified because the baseline import leaves upstream's lockfile untouched, so all three artefacts resolve the same \texttt{libcrux-ml-kem} 0.0.2. The KEM-module diff between the baseline and \texttt{v2.0.0} also adds the HQC arm without altering the existing ones. A companion Criterion target measures ML-DSA-87 key generation, signing, and verification.

The session harness implements two explicitly bounded handshake spans. The initiator span begins by processing the responder's bundle and concludes after serialising the initial \texttt{PreKeySignalMessage}. The responder span begins with that serialised message and a fresh store containing the corresponding private pre-keys. It concludes after successful decryption and insertion of the established \texttt{SessionRecord} into the in-memory \texttt{SessionStore}. This endpoint updates an in-memory store, not durable persistence or a measurement of the \texttt{SessionRecord} encoding. Criterion fixture setup creates or clones each fresh store outside the timed region, ensuring that store preparation is excluded from both spans. Versions 1 and 2 are each measured in last-resort-only and last-resort-plus-one-time modes. The baseline is measured in separately named configurations with and without its optional curve one-time pre-key. These configurations are not considered semantically equivalent to the two KEM modes of the variants. Any retained steady-state encryption or decryption benchmarks are reported separately and not described as handshake latency.

The baseline harness is derived from project commit \texttt{26ff061}. The Version 1 and Version 2 harnesses are derived from their protocol tags and use identical timing boundaries. Exact commands, full harness commits, feature sets, run order, machine identifier, and independent invocation identifiers accompany the archived outputs. For example, the common KEM target is invoked with \texttt{cargo bench -p libsignal-protocol --features "hqc256 mlkem1024" --bench kem}.

Each benchmark-matrix cell is executed in five independent Criterion invocations per machine. Invocations are organised into five complete blocks, each containing every applicable artefact, configuration, and benchmark case once. Case order is cyclically rotated between blocks according to the archived schedule. For each cell, Chapter~\ref{sec:results} reports all five invocation-level medians, their median, and their full range. A relative ordering is classified as stable only when it has the same direction in all five blocks and the ranges of the invocation-level medians do not overlap. Otherwise, the comparison is reported without a stable ordering. Criterion's within-invocation confidence intervals are reported separately and are not pooled across invocations. Each machine's results are analysed independently, with no averaging or statistical comparison across machines.

%Chapter 4.4
\subsection{Size and Storage Measurement}
\label{subsec:size-measurement}

Size evidence differentiates among raw primitive lengths, complete serialised protocol objects, and analytical storage models. A serialised message length is byte-exact for a pinned serialiser only when the mode, plaintext, registration identifier, pre-key identifiers, counters, and field presence are fixed. The size harness therefore employs the same canonical values for every artefact and on every invocation. The common primitive harness reports the lengths of the raw public key, secret key, ciphertext, and shared secret. For key and ciphertext encodings, it records both the tag-free primitive length and the artefact's type-prefixed length. Shared secrets carry no algorithm tag. Secret-key and shared-secret lengths represent local cryptographic material rather than protocol-wire traffic.

For the baseline, Version 1, and Version 2, a dedicated capture serialises the initial \texttt{PreKeySignalMessage} and the first reply under each named configuration using the fixed inputs above. Only the initial message is compared with the lower bounds of Table~\ref{tab:spec-sizes}. The difference includes numeric identifiers, outer and inner version bytes, Protocol Buffers field and length framing, the nested \texttt{SignalMessage}, its ratchet key and counters, the padded encrypted payload, and the eight-byte message authentication code (MAC). The first reply is reported separately. Both values represent libsignal protocol-object bytes. They exclude service envelopes, destination metadata, Sealed Sender, WebSocket or TLS framing, and other transport overhead, and are therefore not measurements of application-level network bandwidth.

\texttt{PreKeyBundle} is an in-memory aggregate and has neither a byte serialiser nor a Protocol Buffers representation in the measured crate. Therefore, no complete serialised bundle length is reported. Each published key is instead measured in its type-prefixed encoding and each signature by its raw length. These values are summed as a \emph{serialised cryptographic bundle-component total}. Client-local record encodings include private keys and session state and are neither bundle wire formats nor server-storage proxies. For Versions 1 and 2, server-held public cryptographic payload per device is modelled as a function of the number \(N\) of one-time KEM pre-keys, consistent with Section~\ref{subsec:theoretical-comparison}. If the baseline's one-time KEM and curve-key pools are not fixed by a single provisioning policy, their counts are parameterised separately. These expressions are analytical payload models rather than direct measurements of a Signal server. They exclude record framing, databases and indexes, at-rest protection, replication, allocator or page overhead, and service metadata.

%Chapter 4.5
\subsection{Threats to Validity}
\label{subsec:threats-validity}

The full threat analysis is given in Appendix~\ref{app:threats-detail}; this section states its substance. Internally, timings measure pinned library integrations, so library maturity and optimisation remain confounded with the choice of KEM, and baseline-to-Version-1 results measure the aggregate transition rather than any single primitive. As constructs, the baseline is a historical pre-SPQR configuration, the optional modes are named configurations rather than equivalents, and handshake latency denotes only the bounded spans of Section~\ref{subsec:perf-measurement}. For conclusions, consumer machines are not controlled rigs, so the five-block protocol reports invocation-level medians with ranges and admits a stable ordering only under the pre-specified non-overlap rule, within a single machine. Externally, one implementation per primitive on two desktop platforms bounds generalisation, and no claim is made about handsets. On security, functional tests do not formally verify the modified handshakes, several bindings and enforcement points remain incomplete or unresolved, and Chapter~\ref{sec:results} must distinguish demonstrated, conditional, unresolved, and unmet properties.
